\begin{definition}
\label{definition-principal-homogeneous-space}
Let $S$ be a scheme.
Let $(G, m)$ be a group scheme over $S$.
Let $X$ be a pseudo $G$-torsor over $S$.
\begin{enumerate}
\item We say $X$ is a {\it principal homogeneous space}
or a {\it $G$-torsor} if there exists a fpqc covering\footnote{This means
that the default type of torsor is a pseudo torsor which is trivial on an
fpqc covering. This is the definition in \cite[Expos\'e IV, 6.5]{SGA3}.
It is a little bit inconvenient for us as we most often work in the fppf
topology.}
$\{S_i \to S\}_{i \in I}$ such that each
$X_{S_i} \to S_i$ has a section (i.e., is a trivial pseudo $G_{S_i}$-torsor).
\item Let $\tau \in \{Zariski, \etale, smooth, syntomic, fppf\}$.
We say $X$ is a {\it $G$-torsor in the $\tau$ topology}, or a
{\it $\tau$ $G$-torsor}, or simply a {\it $\tau$ torsor}
if there exists a $\tau$ covering $\{S_i \to S\}_{i \in I}$
such that each $X_{S_i} \to S_i$ has a section.
\item If $X$ is a $G$-torsor, then we say that it is
{\it quasi-isotrivial} if it is a torsor for the \'etale topology.
\item If $X$ is a $G$-torsor, then we say that it is
{\it locally trivial} if it is a torsor for the Zariski topology.
\end{enumerate}
\end{definition}